\documentclass{subfiles}
\begin{document}
    In this section we provide the reader with the notions of group, ring and field;
        in particular, we describe cyclic groups and finite fields.
        In addition, we show how, given some field, we can create a simple extention.

    Conceptually, the notion of group is very simple:
        we have a set $G$ and a binary operation $+$ defined on $G$,
        defined in such a way the the following properties hold:
        \begin{enumerate}
            \item There exists $e \in G$ such that $e + a = a + e = a$;
                that is, there is a neutral element for $+$.

            \item For any $a \in G$, there exists $-a$ such that $-a + a = a + (-a) = e$,
                i.e., each element has an inverse.

            \item For all $a, b, c \in G$, the equality $a + (b + c) = (a + b) + c$ holds.
        \end{enumerate}
        Additionally,if one has that $a + b = b + a$ for all $a, b \in G$,
        we say that $(G, +)$ is a commutative (or abelian) group.

    Since groups are fundamentally sets, it makes sense that these have subsets,
        or more precisely subgroups. More precisely, if $(G, +)$ is a group and 
        $H \subseteq G$ is non-empty, then $(H, +)$ is a sub-group of $(G, +)$ if $(H, +)$ is a group.
        
    One of the most important results in group theory,
        is the so-called \emph{Lagrange's theorem},
        which roughuly states that, for finite groups, i.e, groups with a finite amount of elements,
        the order of a subgroups (the number of elements in the subgroups) divides the order of the group.

    Let us consider $(G, +)$ a group, and $a \in G$ an element.
        Let $a^{i}$ denote the $i$-th power of $a$.
        We say that a certain integer $i$ is the order of $a$ if $a^{i} = e$ and there is no smaller integer satisfying this property. 
        We denote by $\gen{a}$ the set of all powers of an element $a$.
        If $G = \gen{a}$, then we say that $G$ is a cyclic group.
    
    Although groups are very important objects in math,
        so much so that a whole field of study developed around them,
        oftentimes we need something more powerful.
        Luckly for us, more powerful algebraic structures exist,
        and in this section we consider rings.

    As for groups, the notion of ring is in and of itself, really simple:
        we have a set $R$, and two binary operations $+, \cdot$,
        such that $(R, +)$ forms an abelian group and the following properties hold:
        \begin{enumerate}
            \item For all $a, b, c \in R, a \cdot (b \cdot c) = (a \cdot b) \cdot c$.
            \item For all $a, b, c \in R, a \cdot (b + c) = a \cdot b + a \cdot c$. 
        \end{enumerate}
        Moreover, if there exists $e' \in G$ such that $a \cdot e' = e' \cdot a = a$,
        and for all $a \in G$ there exists $a^{-1}$ such that $a \cdot a^{-1} = a^{-1} \cdot a = e'$,
        we say that $(R, +, \cdot)$ is a field.

    It turns out that most of the results that hold for groups, have a counter part in terms of rings;
        that is, we have the notion of sub-ring, order of a ring and an equivalent Lagrange's theorem.
        But all of these results, are far beyond our interests.
        In what follows, we denote by $\K$ a generic field,
        and by $\gf(p)$ a finite field with $p$ elements, where $p$ is usually a prime.
\end{document}
